\section{Momentum smearing: the basic idea}
\label{sec:basic}
\begin{figure}
\includegraphics[width=0.48\textwidth]{momsmear_basic_principle_14pt.pdf}
\caption{Conventional smearing versus momentum smearing for the example
of a Gaussian wave function in $d=1$ spatial dimensions.
The momentum $k$ shifts the centre of the distribution
in momentum space, resulting in an oscillatory behaviour in position space.}
\label{fig:transform}
\end{figure}

As discussed above, quark smearing within hadronic sources or sinks is essential in lattice simulations
to increase the overlap with the desired physical state, reflecting the fact that hadrons are extended objects,
rather than pointlike. A smearing operator $F$ is diagonal in time, trivial
in spin and acts
on the position and colour indices of quark fields:
\begin{equation}
	(F q)_{\vect{x}} = \sum_{\vect{y}\in (a \mathbb{Z})^d}\, f_{\vect{x}-\vect{y}}\, \GC_{\vect{x}\vect{y}}\, q_{\vect{y}} \label{formwupN}\,,
\end{equation}
where $f$ is a scalar function, $\GC$ is a gauge covariant transporter,
which in the free case will be a unit matrix in colour and position space,
and $d$ is the number of spatial dimensions, usually $d=3$.
Note that the field $q_{\vect{x}}$ is usually periodic in $\vect{x}$ on the lattice, whereas $f_{\vect{x}-\vect{y}}$ need not be periodic in $\vect{x}-\vect{y}$. 
In the free case, the convolution Eq.~\eqref{formwupN} becomes a product in Fourier space
\begin{equation}
	\sum_{\vect{x}\in (a \mathbb{Z})^d} e^{i\vect{p}\cdot\vect{x}}\,(F q)_{\vect{x}} = 
	\tilde{f}(\vect{p})\, \tilde{q}_{\vect{p}}\,.	\label{wupfour}
\end{equation}
For the special case of a Gaussian,
\begin{equation}
\label{eq:smgauss}
	 f_{\vect{x}-\vect{y}} = f_{\vect{0}} \exp\left(-\frac{|\vect{x}-\vect{y}|^2}{2{\sigma}^2}\right)\,,
\end{equation}
the Fourier transformed smearing kernel again is a Gaussian:
\begin{equation}
	\tilde{f}(\vect{p}) \equiv  \sum_{\vect{z}\in (a \mathbb{Z})^d} e^{i\vect{p}\cdot\vect{z}} f_{\vect{z}} = \tilde{f}({\vect{0}}) \exp\left( - \frac{\sigma^2\vect{p}^2}{2} \right)  \, .	\label{kerfour}
\end{equation}
Thus, the smeared quark operator has maximal overlap with a quark at rest, $\vect{p}=\vect{0}$.
Non-zero velocities are suppressed in accordance with the above Gaussian momentum distribution.
Clearly, for hadrons carrying significant spatial momenta, such a smearing may be counterproductive.

Having identified the problem, it is easy to modify the smearing to perform well for moving hadrons.
We aim at a momentum distribution of quarks centred around a finite momentum $\vect{k}$, so we need to shift the smearing kernel in momentum space:
\begin{equation}
	\tilde{f}(\vect{p}) \mapsto  \tilde{f}(\vect{p} - \vect{k}) \label{kerfourshift}\,,
\end{equation}
as illustrated in Fig.~\ref{fig:transform}.
This translates to the replacement in position space
\begin{equation}
	f_{\vect{z}} \mapsto e^{i\vect{k}\cdot\vect{z}} f_{\vect{z}} \ .
\end{equation}
Our modified smearing operator $F_{(\vect{k})}$, where $F_{(\vect{0})}=F$,
can thus be formally expressed as
\begin{align}
	(F_{(\vect{k})} q)_{\vect{x}} &= \sum_{\vect{y}\in (a \mathbb{Z})^d} F_{(\vect{k})\vect{x}\vect{y}} q_{\vect{y}}\nonumber\\&= 
	\sum_{\vect{y}\in (a \mathbb{Z})^d} e^{-i\vect{k}\cdot(\vect{x}-\vect{y})}\, f_{\vect{x}-\vect{y}}\, \GC_{\vect{x}\vect{y}}\, q_{\vect{y}} \label{formnewsm}\ .
\end{align}
The only new ingredient is the additional
phase factor $\exp\left[-i\vect{k}\cdot(\vect{x}-\vect{y})\right]$.
Note that the quark momentum shift
$\vect{k}$ need \emph{not} be a lattice momentum, i.e.\ it is not restricted to discrete values $\vect{k}\in(2\pi/L)\mathbb{Z}^d$. The smearing kernel $f$
and the gauge dependent factor $\GC$ can be taken over from any existing smearing method. For an iterative smearing method,
the extra phase factor can easily be integrated into the elementary smearing step. Below we demonstrate this for the
example of Wuppertal smearing.
In principle, there could be additional effects
like a Lorentz contraction of the wave function. This
was studied, e.g., in Refs.~\cite{Roberts:2012tp,DellaMorte:2012xc},
and we will also address this possibility.

The new, modified smearing operator $F_{(\vect{k})}$ of Eq.~\eqref{formnewsm} inherits important properties
from the smearing operator $F_{(\vect{0})}$ it is based on.
If the unmodified smearing operator is self-adjoint, then so is the modified smearing operator $F_{(\vect{k})}$,
because from $f_{\vect{x}-\vect{y}} = f_{\vect{y}-\vect{x}}^*$ and $\GC_{\vect{x}\vect{y}} = \GC^\dagger_{\vect{y}\vect{x}}$ it follows that
$F_{(\vect{k})\vect{x}\vect{y}} = F^\dagger_{(\vect{k})\vect{y}\vect{x}}$.
The underlying smearing operator $F_{(\vect{0})}$ should transform according to an irreducible representation (irrep) of the cubic group $O_h$.
Usually, this will be the trivial $A_1$ representation
but the smearing operator can also be used to inject angular momentum
or to add non-trivial gluonic excitations~\cite{Lacock:1994qx}.
Obviously, $F_{(\vect{k})}$ with a momentum shift $\vect{k}\neq\vect{0}$
breaks the cubic symmetry. However, when used in conjunction 
with a momentum projection that selects hadrons with a momentum $\vect{p}$, then as long
as $\vect{k}\parallel \vect{p}$, the smearing operator remains within the
$A_1$ representation of the $O_h$ little group corresponding to
the momentum direction $\vect{p}$.

